% !TeX root = ../main.tex

\section*{Homework \#13 [2025-12-02]}

\begin{paracol}{2}
  \begin{problem}
    Write down the expression for the contribution to self-energy by the
    above diagram. Discuss whether this diagram contributes to a finite
    quasiparticle lifetime, i.e., whether it has an imaginary part?
  \end{problem}
  \switchcolumn \centering
  \begin{tikzpicture}
    \begin{feynhand}
      \vertex (a) at (0,0);
      \vertex (b) at (1,0);
      \vertex (c) at (2,0);
      \vertex (d) at (3,0);
      \propag [fermion] (a) to (b);
      \propag [fermion] (b) to (c);
      \propag [fermion] (c) to (d);
      \propag [photon, out = 90, in = 90] (a) to (c);
      \propag [photon, out = 90, in = 90] (b) to (d);
    \end{feynhand}
  \end{tikzpicture}
\end{paracol}
\vspace{-12pt}
\begin{solution}
  For convenience, label the vertices, external Fermion momentum, photon
  momenta, and internal Fermion momenta in the Feynman diagram.
  \begin{center}
    \begin{tikzpicture}[scale = 2]
      \begin{feynhand}
        \vertex (a) at (0,0) node at (a) [below] {$a$};
        \vertex (b) at (1,0) node at (b) [below] {$b$};
        \vertex (c) at (2,0) node at (c) [below] {$c$};
        \vertex (d) at (3,0) node at (d) [below] {$d$};
        \propag [fermion] (a) to [edge label' = $k_1$] (b);
        \propag [fermion] (b) to [edge label' = $k_2$] (c);
        \propag [fermion] (c) to [edge label' = $k_3$] (d);
        \propag [photon, out = 90, in = 90] (a) to [edge label = $q_1$] (c);
        \propag [photon, out = 90, in = 90] (b) to [edge label = $q_2$] (d);
        \propag [anti fermion] (a) --++ (-1,.5) node [left] {$p$};
        \propag [fermion] (d) --++ (1,.5) node [right] {$p$};
      \end{feynhand}
    \end{tikzpicture}
  \end{center}
  List the momentum conservations at the four vertices
  \[
    \begin{cases*}
      p = k_1 + q_1,    & Vertex a,\\
      k_1 = k_2 + q_2,  & Vertex b,\\
      k_2 + q_1 = k_3,  & Vertex c,\\
      k_3 + q_2 = p,    & Vertex d.
    \end{cases*}
  \]
  Thus, we can obtain the internal Fermion momenta
  \[
    k_1 = p - q_1, \quad k_2 = p - q_1 - q_2, \quad k_3 = p - q_2.
  \]
  Applying the QED Feynman rules in Feynman gauge
  \begin{enumext*}[columns = 8]
    \item(2) Vertex factor: $-\iu\upe\gamma^\mu$.
    \item(3) Fermion propagator:
    $\frac{\iu(\slashed k + m)}{k^2 - m^2 + \iu\epsilon}$.
    \item(3) Photon propagator: $-\frac{\iu g_{\mu\nu}}{q^2 + \iu\epsilon}$.
  \end{enumext*}
  THe self-energy $\Sigma(p)$ is defined such that the 1PI diagram contributes
  to $-\iu\Sigma(p)$. Reading the diagram from a to d
  \begin{multline*}
    -\iu\Sigma(p) = \int \frac{\d^4q_1}{(2\pi)^4} \int \frac{\d^4q_2}{(2\pi)^4}
    \cdot
    (-\iu\upe\gamma^\alpha)
    \frac{\iu(\slashed k_1 + m)}{k_1^2 - m^2 + \iu\epsilon}
    \cdot
    (-\iu\upe\gamma^\gamma)
    \frac{\iu(\slashed k_2 + m)}{k_2^2 - m^2 + \iu\epsilon}\\
    \cdot
    (-\iu\upe\gamma^\beta)
    \frac{\iu(\slashed k_3 + m)}{k_3^2 - m^2 + \iu\epsilon}
    \cdot
    (-\iu\upe\gamma^\delta)
    \ab[-\frac{\iu g_{\alpha\beta}}{q_1^2 + \iu\epsilon}]
    \cdot
    \ab[-\frac{\iu g_{\gamma\delta}}{q_1^2 + \iu\epsilon}].
  \end{multline*}
  Simplifying the Expression
  \begin{multline*}
    \Sigma(p) = \iu \upe^4
    \int \frac{\d^4q_1}{(2\pi)^4} \int \frac{\d^4q_2}{(2\pi)^4}
    \frac{1}{q_1^2 + \iu\epsilon} \frac{1}{q_2^2 + \iu\epsilon}
    \gamma^\alpha
    \frac{\slashed p - \slashed q_1 + m}{(p - q_1)^2 - m^2 + \iu\epsilon}\\
    \times \gamma^\gamma
    \frac{\slashed p - \slashed q_1 - \slashed q_2 + m}
      {(p - q_1 - q_2)^2 - m^2 + \iu\epsilon}
    \gamma_\alpha
    \frac{\slashed p - \slashed q_2 + m}
      {(p - q_2)^2 - m^2 + \iu\epsilon}
    \gamma_\gamma
  \end{multline*}
  We obtain the final expression for the crossed self-energy diagram.

  The crossed second-order self-energy diagram does not contribute to a finite
  quasiparticle lifetime at zero temperature: it has no imaginary part for an
  on-shell fermion.
  \begin{enumext}
    \item In vacuum QED, an on-shell electron $p^2 = m^2$ cannot decay into an
    electron plus two photons due to energy-momentum conservation. Any physical
    cut of the diagram would require simultaneous on-shell conditions for two
    massless photons and an intermediate fermion, which is kinematically
    forbidden except for trivial configurations of measure zero in phase space.
    \item In a zero-temperature Fermi liquid, the imaginary part of the
    self-energy on the Fermi surface requires available phase space for
    scattering. For this diagram, putting all three internal fermion lines
    on-shell simultaneously is prohibited by Pauli blocking: the required
    intermediate states involve transitions below the Fermi level, which are
    forbidden at $T=0$. The phase space for such processes vanishes identically.
  \end{enumext}
  Thus, this diagram contributes only to the real part of the self-energy,
  renormalizing the quasiparticle energy and effective mass, but does not induce
  a finite lifetime. Any imaginary part would arise only at finite temperature
  or for off-shell momenta, but for the standard on-shell condition at $T = 0$,
  $\Im \Sigma(p) = 0$.
\end{solution}

\begin{problem}
  Identify each of the Feynman diagrams shown in the figure
  \begin{framed}
  \begin{multline*}
    \tikz [baseline = -.25em]
      \node [fill = red,    minimum width = 4em,
             draw, ellipse, minimum height = 2em] {$\Sigma$};
    \, = \,
    \begin{tikzpicture}[baseline = -.25em]
      \begin{feynhand}
        \vertex (a) at (0,0);
        \vertex (b) at (1.6,0);
        \propag [fermion] (a) to (b);
        \propag [photon, out = 90, in = 90] (a) to (b);
      \end{feynhand}
    \end{tikzpicture} \, + \,
    \begin{tikzpicture}[baseline = -1cm - .25em]
      \begin{feynhand}
        \vertex (a) at (-.5,0);
        \vertex (b) at (.5,0);
        \propag [out = -90, in = -90] (a) to (b)
         coordinate [midway] (c);
        \propag [fermion, out = 90, in = 90] (b) to (a);
        \vertex (d) at (0,-1);
        \propag [photon] ([yshift = -.3cm]c) to (d);
      \end{feynhand}
    \end{tikzpicture} \, + \,
    \begin{tikzpicture}[baseline = -1cm - .25em]
      \begin{feynhand}
        \vertex (a) at (-.5,0);
        \vertex (b) at (.5,0);
        \propag [out = -90, in = -90] (a) to (b)
         coordinate [midway] (c);
        \propag [fermion, out = 90, in = 90] (b) to (a);
        \vertex (d) at (0,-1);
        \propag [photon] ([yshift = -.3cm]c) to (d);
        \vertex (e) at ([xshift = -.8cm]d);
        \vertex (f) at ([xshift = +.8cm]d);
        \propag [fermion] (e) to (f);
        \propag [photon, in = -90, out = -90] (e) to (f);
      \end{feynhand}
    \end{tikzpicture} \, + \,
    \begin{tikzpicture}[baseline = -.25em]
    \begin{feynhand}
      \vertex (a) at (0,0);
      \vertex (b) at (1,0);
      \vertex (c) at (2,0);
      \vertex (d) at (3,0);
      \propag [fermion] (a) to (b);
      \propag [fermion] (b) to (c);
      \propag [fermion] (c) to (d);
      \propag [photon, out = 90, in = 90] (a) to (d);
      \propag [photon, out = 90, in = 90] (b) to (c);
    \end{feynhand}
  \end{tikzpicture}\\
  + \,
  \begin{tikzpicture}[baseline = -2.3cm - .25em]
    \begin{feynhand}
      \vertex (a) at (-.5,0);
      \vertex (b) at (.5,0);
      \propag [out = -90, in = -90] (a) to (b)
       coordinate [midway] (c);
      \propag [fermion, out = 90, in = 90] (b) to (a);
      \vertex (d) at (0,-1);
      \propag [photon] ([yshift = -.3cm]c) to (d);
    \end{feynhand}
    \begin{feynhand}[shift = {(.2,-1.3)}]
      \vertex (a) at (-.5,0);
      \vertex (b) at (.5,0);
      \propag [out = -90, in = -90] (a) to (b)
       coordinate [midway] (c);
      \propag [fermion, out = 90, in = 90] (b) to (a);
      \vertex (d) at (0,-1);
      \propag [photon] ([yshift = -.3cm]c) to (d);
    \end{feynhand}
  \end{tikzpicture} \, + \,
    \begin{tikzpicture}[baseline = -1cm - .25em]
      \begin{feynhand}
        \vertex (a) at (-.5,0);
        \vertex (b) at (.5,0);
        \propag [out = -90, in = -90] (a) to (b)
         coordinate [midway] (c);
        \propag [anti fermion, out = 90, in = 90] (b) to (a);
        \vertex (d) at (0,-1);
        \propag [photon] ([yshift = -.3cm]c) to (d);
        \propag [photon, in = 180, out = 180] (a) to (0,.5);
        \propag [photon, in = 0, out = 0] (0,.5) to (b);
      \end{feynhand}
    \end{tikzpicture} \, + \,
    \begin{tikzpicture}
      \begin{feynhand}
        \vertex (a) at (0,0);
        \vertex (b) at (1,0);
        \vertex (c) at (2,0);
        \vertex (d) at (3,0);
        \propag [fermion] (a) to (b);
        \propag [fermion] (b) to (c);
        \propag [fermion] (c) to (d);
        \propag [photon, out = 90, in = 90] (a) to (c);
        \propag [photon, out = 90, in = 90] (b) to (d);
      \end{feynhand}
    \end{tikzpicture} \, + \,
    \begin{tikzpicture}[baseline = -1.5cm - .25em]
      \begin{feynhand}
        \vertex (a) at (-1,0);
        \vertex (b) at (1,0);
        \propag [out = -90, in = -90] (a) to (b)
         coordinate [midway] (c);
        \propag [fermion, out = 90, in = 90] (b) to (a);
        \vertex (c) at (-.6,-.5);
        \vertex (d) at (.6,-.5);
        \vertex (e) at ([yshift = -1cm]c);
        \vertex (f) at ([yshift = -1cm]d);
        \propag [photon] (c) to (e);
        \propag [photon] (d) to (f);
        \propag [fermion] (e) to (f);
      \end{feynhand}
    \end{tikzpicture}
  \end{multline*}
  \end{framed}
  with its corresponding skeleton diagram (where there are just four skeleton
  diagrams); Illustrate why it corresponds to that skeleton diagram.
\end{problem}
\begin{solution}
  The non-skeleton diagram represents the self-energy $\Sigma$ using bare
  propagators $G_0$, including all diagrams up to second order, and those that
  are reducible: self-energy corrections on internal lines.
  \begin{framed}
  \centering
  To distinguish the two figures, double lines are used to describe the full
  propagator $G$, contrary to $G_0$.
  \[
    \tikz [baseline = -.25em]
      \node [fill = red,    minimum width = 4em,
             draw, ellipse, minimum height = 2em] {$\Sigma$};
    \, = \,
    \begin{tikzpicture}[baseline = -.25em]
      \begin{feynhand}
        \vertex (a) at (0,0);
        \vertex (b) at (1.6,0);
        \propag [fermion] ([yshift = +.6pt]a) to ([yshift = +.6pt]b);
        \propag [fermion] ([yshift = -.6pt]a) to ([yshift = -.6pt]b);
        \propag [photon, out = 90, in = 90] (a) to (b);
      \end{feynhand}
    \end{tikzpicture} \, + \,
    \begin{tikzpicture}[baseline = -1cm - .25em]
      \begin{feynhand}
        \vertex (a) at (-.5,0);
        \vertex (b) at (.5,0);
        \propag [out = -90, in = -90]
                ([xshift = +.6pt]a) to ([xshift = -.6pt]b);
        \propag [out = -90, in = -90] ([xshift = -.6pt]a) to ([xshift = +.6pt]b)
          coordinate [midway] (c);
        \propag [fermion, out = 90, in = 90]
                ([xshift = -.6pt]b) to ([xshift = +.6pt]a);
        \propag [fermion, out = 90, in = 90]
                ([xshift = +.6pt]b) to ([xshift = -.6pt]a);
        \vertex (d) at (0,-1);
        \propag [photon] ([yshift = -.3cm]c) to (d);
      \end{feynhand}
    \end{tikzpicture} \, + \,
    \begin{tikzpicture}
      \begin{feynhand}
        \vertex (a) at (0,0);
        \vertex (b) at (1,0);
        \vertex (c) at (2,0);
        \vertex (d) at (3,0);
        \propag [fermion] ([yshift = +.6pt]a) to ([yshift = +.6pt]b);
        \propag [fermion] ([yshift = -.6pt]a) to ([yshift = -.6pt]b);
        \propag [fermion] ([yshift = +.6pt]b) to ([yshift = +.6pt]c);
        \propag [fermion] ([yshift = -.6pt]b) to ([yshift = -.6pt]c);
        \propag [fermion] ([yshift = +.6pt]c) to ([yshift = +.6pt]d);
        \propag [fermion] ([yshift = -.6pt]c) to ([yshift = -.6pt]d);
        \propag [photon, out = 90, in = 90] (a) to (c);
        \propag [photon, out = 90, in = 90] (b) to (d);
      \end{feynhand}
    \end{tikzpicture} \, + \,
    \begin{tikzpicture}[baseline = -1.5cm - .25em]
      \begin{feynhand}
        \vertex (a) at (-1,0);
        \vertex (b) at (1,0);
        \propag [out = -90, in = -90]
                ([xshift = +.6pt]a) to ([xshift = -.6pt]b);
        \propag [out = -90, in = -90] ([xshift = -.6pt]a) to ([xshift = +.6pt]b)
         coordinate [midway] (c);
        \propag [fermion, out = 90, in = 90]
                ([xshift = +.6pt]b) to ([xshift = -.6pt]a);
        \propag [fermion, out = 90, in = 90]
                ([xshift = -.6pt]b) to ([xshift = +.6pt]a);
        \vertex (c) at (-.6,-.5);
        \vertex (d) at (.6,-.5);
        \vertex (e) at ([yshift = -1cm]c);
        \vertex (f) at ([yshift = -1cm]d);
        \propag [photon] (c) to (e);
        \propag [photon] (d) to (f);
        \propag [fermion] ([yshift = +.6pt]e) to ([yshift = +.6pt]f);
        \propag [fermion] ([yshift = -.6pt]e) to ([yshift = -.6pt]f);
      \end{feynhand}
    \end{tikzpicture}
  \]
  \end{framed}
  In contrast, the skeleton expansion defines $\Sigma$ as a functional of the
  fully interacting Green's function $G$. By using $G$, the expansion includes
  only the topologically irreducible diagrams.

  \paragraph{Identification for the 4 skeleton diagrams}
  \begin{enumext}
    \item \textbf{Fock (Exchange) Diagram} corresponding to the following 2
    terms in the non-skeleton diagram.
    \begin{enumext}[wrap-label = \underline{Term #1}:, label = \arabic*,
                    list-indent = 0pt]
      \setcounter{enumXii}{0}
      \item \textsc{Direct Correspondence.} This is the bare 1st-order
      Fock (Exchange) diagram, which defines the basic topology of the skeleton;
      and the
      \setcounter{enumXii}{6}
      \item \textsc{Reducible Diagram (by deduction).} This diagram is absorbed
      into this skeleton diagram. While its exact internal
      structure is complex, it is a second-order diagram that is topologically
      equivalent to the Fock diagram once its internal lines are treated as
      fully interacting $G$ lines, thereby implicitly summing this type of
      correction.
    \end{enumext}
    \item \textbf{Hartree (Tadpole) Diagram} corresponding to the following 4
    terms in the non-skeleton diagram.
    \begin{enumext}[wrap-label = \underline{Term #1}:, label = \arabic*,
                    list-indent = 0pt]
      \setcounter{enumXii}{1}
      \item \textsc{Direct Correspondence.} This is the bare 1st-order Hartree
      (Tadpole) diagram, which defines the basic topology of the skeleton.
      \setcounter{enumXii}{2}
      \item \textsc{Reducible Diagram.} This is obtained by inserting a
      1st-order self-energy correction (a Fock-like semi-circle) into the
      internal Green's function line of the Tadpole diagram. This correction is
      automatically included when the internal line in this skeleton diagram is
      defined as $G$.
      \setcounter{enumXii}{4}
      \item \textsc{Reducible Diagram.} This is a more complex version of a
      correction to the Tadpole diagram's internal line. It is absorbed into
      this skeleton diagram.
      \setcounter{enumXii}{5}
      \item \textsc{Reducible Diagram.} This is obtained by inserting a
      1st-order self-energy correction (a Hartree-like tadpole) into the
      internal Green's function line of the Tadpole diagram. It is absorbed into
      this skeleton diagram.
    \end{enumext}
    \item \textbf{Second-Order Sunset Diagram} corresponding to the 4th term in
    the non-skeleton diagram: \textsc{Direct Correspondence}. This is the bare
    second-order diagram with the Sunset (or double-semi-circle) topology. It is
    topologically irreducible that cannot be reduced to a 1st-order self-energy
    by cutting a single internal line, and it is retained as a distinct
    skeleton.
    \item \textbf{Second-Order Crossed/Double-Tadpole Diagram} corresponding to
    the 8th term in the non-skeleton diagram: \textsc{Direct Correspondence}.
    This is the bare second-order diagram with the Double-Tadpole (or hourglass)
    topology. It is also topologically irreducible and is retained as a distinct
    skeleton.
  \end{enumext}
  \paragraph{Illustration for the correspondence}
  The reducible diagrams (Terms 3, 5, 6, and 7) in the non-skeleton diagram are
  the four ``self-energy corrections of the internal Green's functions''. By
  using the fully interacting propagator G in the skeleton expansion, these
  corrections are effectively summed up to all orders and accounted for by the
  first two skeletons. In order words, reorganizing the perturbation series into
  a self-consistent formulation.
\end{solution}